% GENERATED FILE. Edit canonical lessons, then run npm run generate:books.
% canonical-source-hash: fnv1a64:b4e09f914b6d4651
% canonical-lessons: GU-C215-tharvu, GU-C215-him, GU-C215-os, GU-C215-gadh, GU-C215-huphalu

\chapter{A cold foggy morning}
\label{ch:gu-a-cold-foggy-morning}
\hlchaptermodality{215}

\begin{chapteropening}
\textbf{By the end of this chapter:} \emph{I can say that it will freeze and that the morning will be foggy.}

Complete the last lesson of chapter 215: I can say that it will freeze and that the morning will be foggy.
\end{chapteropening}

\section[\texorpdfstring{\gu{ઠરવું} (ṭharvũ)}{ઠરવું (ṭharvũ)}]{\gu{ઠરવું} (ṭharvũ) — to freeze, to cool down}

\label{lesson:GU-C215-tharvu}



\begin{quote}
Before the new one: say the Gujarati for to drip, then the Gujarati for to get soaked.
\end{quote}

\subsection*{You'll want to know: \gu{ઠરવું}}
\textbf{\gu{ઠરવું}} — \emph{ṭharvũ} — \textquotedblleft{}to freeze, to cool down\textquotedblright{}.

\begin{quote}
\textbf{\gu{રાત્રે} \gu{પાણી} \gu{ઠરશે}.} — \emph{rātre pāṇī ṭharśe} — \textquotedblleft{}The water will freeze at night.\textquotedblright{}
\end{quote}

\begin{cousinweb}
\textbf{\gu{ઠરવું}} (\emph{ṭharvũ}) is a tadbhava. It comes from \emph{*stharati}, \textquotedblleft{}is firm\textquotedblright{}, a reconstructed form; CDIAL glosses the Gujarati verb \textquotedblleft{}to become fixed, be settled, freeze, congeal\textquotedblright{} (CDIAL 13742).
\end{cousinweb}

\begin{grammarlens}[title={What you've built}]
One more word for everyday life at A2.
\end{grammarlens}

\subsection*{Your turn}
\begin{itemize}
\raggedright
  \item \emph{Say it:} \emph{ṭharvũ}
  \item \emph{Say it:} \emph{ṭharvũ}, once more
  \item \emph{Say it:} \emph{rātre pāṇī ṭharśe} — The water will freeze at night.
  \item \emph{Recall:} say the Gujarati for to drip, then the Gujarati for to get soaked, then say \emph{ṭharvũ} again
\end{itemize}

\begin{tcolorbox}[breakable,colback=teal!4,colframe=teal!35!black,title={Before you move on}]
What does \gu{ઠરવું} mean? (\textquotedblleft{}To freeze, to cool down\textquotedblright{}.) Say it once more.
\end{tcolorbox}
\section[\texorpdfstring{\gu{હિમ} (him)}{હિમ (him)}]{\gu{હિમ} (him) — frost, snow}

\label{lesson:GU-C215-him}



\begin{quote}
Before the new one: say the Gujarati for to get soaked, then the Gujarati for to freeze.
\end{quote}

\subsection*{You'll want to know: \gu{હિમ}}
\textbf{\gu{હિમ}} — \emph{him} — \textquotedblleft{}frost, snow\textquotedblright{}.

\begin{quote}
\textbf{\gu{આવતીકાલે} \gu{સવારે} \gu{હિમ} \gu{પડશે}.} — \emph{āvtīkāle savāre him paḍśe} — \textquotedblleft{}There will be frost tomorrow morning.\textquotedblright{}
\end{quote}

A neuter noun.

\begin{cousinweb}
\textbf{\gu{હિમ}} (\emph{him}) is Sanskrit \emph{hima}, \textquotedblleft{}cold, frost, snow\textquotedblright{}, taken over almost whole, a tatsama. CDIAL lists Gujarati \emph{him}, \textquotedblleft{}hoarfrost, dew\textquotedblright{} (CDIAL 14096).
\end{cousinweb}

\begin{grammarlens}[title={What you've built}]
One more word for everyday life at A2.
\end{grammarlens}

\subsection*{Your turn}
\begin{itemize}
\raggedright
  \item \emph{Say it:} \emph{him}
  \item \emph{Say it:} \emph{him}, once more
  \item \emph{Say it:} \emph{āvtīkāle savāre him paḍśe} — There will be frost tomorrow morning.
  \item \emph{Recall:} say the Gujarati for to get soaked, then the Gujarati for to freeze, then say \emph{him} again
\end{itemize}

\begin{tcolorbox}[breakable,colback=teal!4,colframe=teal!35!black,title={Before you move on}]
What does \gu{હિમ} mean? (\textquotedblleft{}Frost, snow\textquotedblright{}.) Say it once more.
\end{tcolorbox}
\section[\texorpdfstring{\gu{ઓસ} (os)}{ઓસ (os)}]{\gu{ઓસ} (os) — dew}

\label{lesson:GU-C215-os}



\begin{quote}
Before the new one: say the Gujarati for to freeze, then the Gujarati for frost.
\end{quote}

\subsection*{You'll want to know: \gu{ઓસ}}
\textbf{\gu{ઓસ}} — \emph{os} — \textquotedblleft{}dew\textquotedblright{}.

\begin{quote}
\textbf{\gu{સવારે} \gu{ઓસ} \gu{પડશે}.} — \emph{savāre os paḍśe} — \textquotedblleft{}Dew will fall in the morning.\textquotedblright{}
\end{quote}

A feminine noun.

\begin{cousinweb}
\textbf{\gu{ઓસ}} (\emph{os}) is a tadbhava, a word inherited from Sanskrit and reshaped by sound change. It comes from \emph{avaśyā}, \textquotedblleft{}hoarfrost, dew\textquotedblright{} (CDIAL 855).
\end{cousinweb}

\begin{grammarlens}[title={What you've built}]
One more word for everyday life at A2.
\end{grammarlens}

\subsection*{Your turn}
\begin{itemize}
\raggedright
  \item \emph{Say it:} \emph{os}
  \item \emph{Say it:} \emph{os}, once more
  \item \emph{Say it:} \emph{savāre os paḍśe} — Dew will fall in the morning.
  \item \emph{Recall:} say the Gujarati for to freeze, then the Gujarati for frost, then say \emph{os} again
\end{itemize}

\begin{tcolorbox}[breakable,colback=teal!4,colframe=teal!35!black,title={Before you move on}]
What does \gu{ઓસ} mean? (\textquotedblleft{}Dew\textquotedblright{}.) Say it once more.
\end{tcolorbox}
\section[\texorpdfstring{\gu{ગાઢ} (gāḍh)}{ગાઢ (gāḍh)}]{\gu{ગાઢ} (gāḍh) — thick, dense}

\label{lesson:GU-C215-gadh}



\begin{quote}
Before the new one: say the Gujarati for frost, then the Gujarati for dew.
\end{quote}

\subsection*{You'll want to know: \gu{ગાઢ}}
\textbf{\gu{ગાઢ}} — \emph{gāḍh} — \textquotedblleft{}thick, dense\textquotedblright{}.

\begin{quote}
\textbf{\gu{કાલે} \gu{સવારે} \gu{ગાઢ} \gu{ધુમ્મસ} \gu{હશે}.} — \emph{kāle savāre gāḍh dhummas haśe} — \textquotedblleft{}There will be thick fog tomorrow morning.\textquotedblright{}
\end{quote}

It does not change for gender.

\begin{cousinweb}
\textbf{\gu{ગાઢ}} (\emph{gāḍh}) comes from Sanskrit \emph{gāḍha}, \textquotedblleft{}pressed together, thick, firm\textquotedblright{}. CDIAL lists Gujarati \emph{gāḍh}, \textquotedblleft{}thick, sound (of sleep)\textquotedblright{} (CDIAL 4118).
\end{cousinweb}

\begin{grammarlens}[title={What you've built}]
One more word for everyday life at A2.
\end{grammarlens}

\subsection*{Your turn}
\begin{itemize}
\raggedright
  \item \emph{Say it:} \emph{gāḍh}
  \item \emph{Say it:} \emph{gāḍh}, once more
  \item \emph{Say it:} \emph{kāle savāre gāḍh dhummas haśe} — There will be thick fog tomorrow morning.
  \item \emph{Recall:} say the Gujarati for frost, then the Gujarati for dew, then say \emph{gāḍh} again
\end{itemize}

\begin{tcolorbox}[breakable,colback=teal!4,colframe=teal!35!black,title={Before you move on}]
What does \gu{ગાઢ} mean? (\textquotedblleft{}Thick, dense\textquotedblright{}.) Say it once more.
\end{tcolorbox}
\section[\texorpdfstring{\gu{હૂંફાળું} (hū̃phāḷũ)}{હૂંફાળું (hū̃phāḷũ)}]{\gu{હૂંફાળું} (hū̃phāḷũ) — lukewarm, pleasantly warm}

\label{lesson:GU-C215-huphalu}



\begin{quote}
Before the new one: say the Gujarati for dew, then the Gujarati for thick.
\end{quote}

\subsection*{You'll want to know: \gu{હૂંફાળું}}
\textbf{\gu{હૂંફાળું}} — \emph{hū̃phāḷũ} — \textquotedblleft{}lukewarm, pleasantly warm\textquotedblright{}.

\begin{quote}
\textbf{\gu{હું} \gu{હૂંફાળું} \gu{પાણી} \gu{પીશ}.} — \emph{hũ hū̃phāḷũ pāṇī pīś} — \textquotedblleft{}I will drink warm water.\textquotedblright{}
\end{quote}

\begin{cousinweb}
\textbf{\gu{હૂંફાળું}} (\emph{hū̃phāḷũ}) is built on \emph{hū̃ph}, \textquotedblleft{}warmth\textquotedblright{}, a tadbhava: CDIAL lists Gujarati \emph{hũph}, \textquotedblleft{}genial warmth\textquotedblright{}, under Sanskrit \emph{ūṣman}, \textquotedblleft{}heat, steam\textquotedblright{} (CDIAL 2441). The next entry, \emph{ūṣmāyate}, \textquotedblleft{}emits heat\textquotedblright{}, holds the close Gujarati \emph{ũphāylũ}, \textquotedblleft{}lukewarm, tepid\textquotedblright{} (CDIAL 2442).
\end{cousinweb}

\begin{grammarlens}[title={What you've built}]
One more word for everyday life at A2.
\end{grammarlens}

\subsection*{Your turn}
\begin{itemize}
\raggedright
  \item \emph{Say it:} \emph{hū̃phāḷũ}
  \item \emph{Say it:} \emph{hū̃phāḷũ}, once more
  \item \emph{Say it:} \emph{hũ hū̃phāḷũ pāṇī pīś} — I will drink warm water.
  \item \emph{Recall:} say the Gujarati for dew, then the Gujarati for thick, then say \emph{hū̃phāḷũ} again
\end{itemize}

\begin{tcolorbox}[breakable,colback=teal!4,colframe=teal!35!black,title={Before you move on}]
What does \gu{હૂંફાળું} mean? (\textquotedblleft{}Lukewarm, pleasantly warm\textquotedblright{}.) Say it once more.
\end{tcolorbox}
